% SPI chapter. Section order is the per-peripheral template enforced by
% python/check_intro_names.py. Vector numbers are frozen by the generator
% (python/generate.py, _mcuMpIrqVectors). Pin lists are generated, not
% written here.
\section{Overview}

The Serial Peripheral Interface (\peripheral{SPIx}) is a full-duplex synchronous serial interface. A bus has three shared wires (clock, master out, master in) and one chip select per slave. \peripheral{SPI0} is the boot-flash master and carries the memory-mapped flash window; \peripheral{SPI1}, in configurations that include it, is a general-purpose port with both master and slave modes.

\begin{itemize}
	\item Master mode on both ports; slave mode on \peripheral{SPI1}
	\item 8-, 16- or 32-bit frames (\bitfield{SPIDL}), LSB- or MSB-first (\bitfield{SPIMSB})
	\item SPI modes 0 to 3 as master (\bitfield{SPICPOL}, \bitfield{SPICPHA}); modes 1 and 3 as slave
	\item Optional byte swap on transmit and receive (\bitfield{SPITXSB}, \bitfield{SPIRXSB})
	\item Baud divider from SMCLK (\bitfield{SPIBR})
	\item One-deep transmit queue for back-to-back frames without dead time
	\item Two interrupts per port: transfer complete, transmit register empty
	\item Flash extended memory on \peripheral{SPI0}: memory-mapped reads and execute-in-place from the boot flash for hart 0 (\bitfield{SPIFEN}, \register{SPI0FOS})
\end{itemize}

\section{Block diagram}

Figure \ref{fig:spi-connection-diagram} shows how the two ports are wired in this configuration: \peripheral{SPI0} to the external boot flash on fixed pads, and \peripheral{SPI1} to the pool of relocatable pins. The flash chip select is a GPIO pad at reset, driven by the boot ROM while it loads the image, and bit 0 of \register{P0SEL} hands the pad driver to \peripheral{SPI0}. Only hart 0 reaches the flash window, and it stalls while \peripheral{SPI0} reads on its behalf (Section \ref{ss:spi-flash-xip}).

\begin{figure}[htbp]
	\centering
	% Generated by LatexUserGuide.GenerateSpiFlashDiagram (make chip).
	\input{include/SpiFlashDiagram.tex}
	\caption{SPI connection diagram: \peripheral{SPI0} to the boot flash on fixed pads and \peripheral{SPI1}, drawn as master, to an external device.}
	\label{fig:spi-connection-diagram}
\end{figure}

\section{Functional description}

\subsection{Transfer}

The bus is idle until the master asserts (drives low) the chip select of one slave. The master then drives the clock for one frame of 8, 16 or 32 bits; on each clock transition the master and the slave alternately update and latch the two data wires, so each frame moves one word in each direction at once. After a frame the master either sends another frame with the chip select still low, which makes the frames one packet, or deasserts the chip select. The slave must have a word ready for the first frame of a packet, because the master clocks it out whether or not the slave has data to send.

The order of update and latch, and the idle level of the clock, are set by clock polarity (CPOL) and clock phase (CPHA), together called the SPI mode (Table \ref{t:SPIMODE-key}). Figure \ref{fig:spi-timing-diagram} draws all four modes for one 8-bit frame. Slave mode supports modes 1 and 3 only (CPHA = 1), because with CPHA = 0 the slave would have to present its first bit before the first clock edge, before it has seen the chip select settle.

\begin{figure}[htpb]
	% Generated by LatexUserGuide.GenerateSpiTimingDiagram (make chip).
	\input{include/SpiTimingDiagram.tex}
	\caption{SPI timing for the four modes, one 8-bit frame.}
	\label{fig:spi-timing-diagram}
\end{figure}

\begin{table}[htb]
	\centering
	\caption{SPI mode key.}
	\label{t:SPIMODE-key}
	\begin{tabular}{ c c c c c c }
	\textbf{Mode} & \textbf{CPOL} & \textbf{CPHA} & \textbf{Clock idle} & \textbf{Data update edge} & \textbf{Data latch edge} \\ \hline
	0 & 0 & 0 & Low & Trailing (falling) & Leading (rising) \\
	\rowcolor{tablehighlightcolor}
	1 & 0 & 1 & Low & Leading (rising) & Trailing (falling) \\
	2 & 1 & 0 & High & Trailing (rising) & Leading (falling) \\
	\rowcolor{tablehighlightcolor}
	3 & 1 & 1 & High & Leading (falling) & Trailing (rising) \\
	\hline
	\end{tabular}
\end{table}

\subsection{Bit and byte order}

\bitfield{SPIMSB} selects MSB-first (1) or LSB-first (0) within a frame. For 16- and 32-bit frames, \bitfield{SPITXSB} and \bitfield{SPIRXSB} swap the byte order of the transmitted and received words, so a word exchanged with a device of the other endianness needs no software swap. The bit order is applied before the byte swap: a 16-bit MSB-first frame with \bitfield{SPITXSB} set sends bits 7 to 0 and then bits 15 to 8. Figures \ref{fig:spi-byte-ordering} and \ref{fig:spi-bit-ordering} show both.

\begin{figure}[htpb]
	% Generated by LatexUserGuide.GenerateSpiByteOrderingDiagram (make chip).
	\input{include/SpiByteOrderingDiagram.tex}
	\caption{Byte order of a multi-byte frame with and without the swap.}
	\label{fig:spi-byte-ordering}
\end{figure}

\begin{figure}[htpb]
	% Generated by LatexUserGuide.GenerateSpiBitOrderingDiagram (make chip).
	\input{include/SpiBitOrderingDiagram.tex}
	\caption{Bit order of a 16-bit frame.}
	\label{fig:spi-bit-ordering}
\end{figure}

\subsection{Baud clock}

In master mode the clock frequency is
\begin{equation*}
f_\textrm{SCK} = \frac{f_\textrm{SMCLK}}{2 \times (1 + \bitfield{SPIBR})}
\end{equation*}
so the fastest clock is half SMCLK. The port's serial engine runs on SMCLK and its register interface on MCLK, so the registers stay readable with SMCLK stopped.

\subsection{Transmit queue}

A write to \register{SPIxTX} queues one word. In master mode with no frame in progress the frame starts at once and \bitfield{SPITEIF} is set immediately; with a frame in progress the queued word starts when that frame ends, and \bitfield{SPITEIF} is set then. In slave mode the last word written is loaded when the master's first clock edge arrives, and \bitfield{SPITEIF} is set at that edge. The queue holds one word, so a second write during one frame overwrites the first; software waits for \bitfield{SPITEIF} = 1 or \bitfield{SPIBUSY} = 0 before each write.

\subsection{Status flags}

\bitfield{SPITCIF} is set when a frame completes and the received word is in \register{SPIxRX}; it is cleared by writing 1 or by reading \register{SPIxRX}. \bitfield{SPITEIF} is set as described above and cleared by writing 1. \bitfield{SPIBUSY} is 1 while the master is clocking a frame, and in slave mode while the chip select is asserted; it can read 0 between frames of one packet when no word was queued. Flash extended memory reads set none of the flags, because they are completed in hardware without software involvement.

\subsection{Flash extended memory} \label{ss:spi-flash-xip}

With \bitfield{SPIFEN} set, hart 0's loads and instruction fetches at or above \texttt{\FlashBaseAddress} become flash reads: \peripheral{SPI0} drives \pin{CS\_FLASH}, sends the read command and a 24-bit address, and stalls hart 0 until the word returns. The window is decoded on hart 0's private bus, not through the arbiter, so the other harts have no path to the flash and read zeros there without stalling. The flash address is
\begin{equation*}
\text{Flash address} = (\text{CPU address} + \register{SPI0FOS}) \bmod 2^{24}
\end{equation*}
so \register{SPI0FOS} relocates the flash contents within the window without moving data; to place flash address 0 at \texttt{\FlashBaseAddress} itself, \register{SPI0FOS} holds the two's complement of the window base modulo $2^{24}$. Each 32-bit read costs a command, an address and four data bytes on the serial clock, so the window suits read-only data, lookup tables and code that is not performance-critical.

\subsection{Boot use of \peripheral{SPI0}}

The boot ROM uses \peripheral{SPI0} to load the application from the flash (Section \ref{s:startup}). During that sequence the ROM drives \pin{CS\_FLASH} as a GPIO output and the flash drives \pin{MISO0}; at the end of the sequence the ROM puts the flash into deep power-down. After boot the port is free for ordinary transfers and for the extended memory window.

\section{Interrupts and events}

Each flag has an enable in \register{SPIxCR} (\bitfield{SPITCIE}, \bitfield{SPITEIE}); the interrupt is a level asserted while both are set. Table \ref{t:spi-irq} lists the vectors.

\begin{table}[htb]
	\centering
	\caption{\peripheral{SPIx} interrupt vectors.}
	\label{t:spi-irq}
	\begin{tabular}{ l l p{4.2cm} l p{4.4cm} }
	\textbf{Vector} & \textbf{Port} & \textbf{Condition} & \textbf{Set by} & \textbf{Cleared by} \\ \hline
	9 & \peripheral{SPI0} & Frame complete & \bitfield{SPITCIF} & Writing 1, or reading \register{SPIxRX} \\
	\rowcolor{tablehighlightcolor}
	10 & \peripheral{SPI0} & Transmit register empty & \bitfield{SPITEIF} & Writing 1 \\
	11 & \peripheral{SPI1} & Frame complete & \bitfield{SPITCIF} & Writing 1, or reading \register{SPIxRX} \\
	\rowcolor{tablehighlightcolor}
	12 & \peripheral{SPI1} & Transmit register empty & \bitfield{SPITEIF} & Writing 1 \\
	\hline
	\end{tabular}
\end{table}

\section{Programming}

\subsection{Initialisation as a master}

\begin{enumerate}
	\item Select the SMCLK source in \register{SYSCLKCR} (Section \ref{ss:clocks}).
	\item Set the \register{PxSEL} bits of the \pin{SCKx}, \pin{MOSIx} and \pin{MISOx} pins, and their \register{PxAFS} fields if the port is relocated.
	\item Configure each slave's chip-select pin as a GPIO output driven high.
	\item Write \register{SPIxCR}.\bitfield{SPIBR} for the clock rate.
	\item Write \register{SPIxCR}.\bitfield{SPICPOL} and \register{SPIxCR}.\bitfield{SPICPHA} for the mode.
	\item Write \register{SPIxCR}.\bitfield{SPIDL} and \register{SPIxCR}.\bitfield{SPIMSB} for the frame format.
	\item Clear \register{SPIxCR}.\bitfield{SPISM}.
	\item Set \register{SPIxCR}.\bitfield{SPIEN}.
\end{enumerate}

\subsection{Master transfer of one packet}

\begin{enumerate}
	\item Write 1 to every flag in \register{SPIxSR}.
	\item Drive the slave's chip-select pin low.
	\item Wait for \bitfield{SPITEIF} = 1 or \bitfield{SPIBUSY} = 0, then write the next word to \register{SPIxTX}.
	\item Wait for \bitfield{SPITCIF} = 1, then read \register{SPIxRX}; the read clears the flag.
	\item Repeat from step 3 for each further frame.
	\item Drive the chip-select pin high.
\end{enumerate}

\subsection{Initialisation as a slave (\peripheral{SPI1})}

\begin{enumerate}
	\item Select the SMCLK source in \register{SYSCLKCR}.
	\item Set the \register{PxSEL} bits of the \pin{SCK1}, \pin{MOSI1} and \pin{MISO1} pins.
	\item Configure the \pin{CS1} pin as a high-impedance GPIO input; the port never drives it.
	\item Write \register{SPI1CR}.\bitfield{SPICPHA} = 1 and \register{SPI1CR}.\bitfield{SPICPOL} for the mode.
	\item Write \register{SPI1CR}.\bitfield{SPIDL} and \register{SPI1CR}.\bitfield{SPIMSB}.
	\item Set \register{SPI1CR}.\bitfield{SPISM}.
	\item Set \register{SPI1CR}.\bitfield{SPIEN}.
	\item Write the first word to \register{SPI1TX}.
\end{enumerate}

In slave mode \bitfield{SPITEIF} at the start of each frame asks for the next word, and \bitfield{SPITCIF} at the end delivers the received one in \register{SPI1RX}; the received word must be read before the next frame ends, because the next frame overwrites it.

\subsection{Enabling the flash window}

\begin{enumerate}
	\item Initialise \peripheral{SPI0} as a master in mode 3 with 8-bit frames at a rate within the flash's specification.
	\item Configure \pin{CS\_FLASH} as a GPIO output driven high.
	\item Send the flash's release-from-deep-power-down command (often \texttt{0xAB}), because the boot ROM leaves the flash in deep power-down.
	\item Send any vendor-specific configuration commands.
	\item Write \register{SPI0FOS}.
	\item Set the \register{PxSEL} bit of \pin{CS\_FLASH}, which hands the pin to \peripheral{SPI0}.
	\item Set \register{SPI0CR}.\bitfield{SPIFEN}.
\end{enumerate}

\section{Usage notes}

\paragraph{Select the SMCLK source before the first transfer.} The boot ROM may leave SMCLK on LFXT (32.768 kHz), where transfers appear to hang.

\paragraph{Never read \register{SPIxRX} speculatively from a hart that does not own the transfer.} The read clears \bitfield{SPITCIF} on any hart's access and consumes the owner's transfer-complete event; when ownership is unclear, clear flags by writing 1s to \register{SPIxSR}.

\paragraph{Serialize whole transfers between harts.} The arbiter makes each register access atomic; frames from two harts interleave on the wire and race for \register{SPIxRX} (Section \ref{ss:multicore-sync}).

\paragraph{Wait for the queue before each write.} A second write to \register{SPIxTX} within one frame overwrites the queued word.

\paragraph{Keep other \peripheral{SPI0} slaves deselected across reset.} The boot ROM clocks the flash while every GPIO is a high-impedance input, so a slave whose chip select can float low drives \pin{MISO0} against the flash; pull the chip-select lines high externally or keep those slaves on \peripheral{SPI1}.

\paragraph{Use the flash window from hart 0 only.} Other harts read zeros there; a program that must run on a channel hart is staged into shared RAM or the hart's TCM first.

\paragraph{Wake the flash before enabling the window.} The window issues read commands only; a flash left in deep power-down returns undefined data.
